\documentclass[../../main.tex]{subfiles}
\graphicspath{{\subfix{../../image/}}}

\begin{document}

\section{迹}\label{section:ch08-变分法}

下面我们讨论在 $\partial U$ 为 $C^{1}$ 的假设下，能否为函数 $u\in W^{1,p}(U)$ 沿 $\partial U$ 赋予“边界值”。若 $u\in C(\overline{U})$，则 $u$ 显然在通常意义下在 $\partial U$ 上有值。问题在于，典型的 $u\in W^{1,p}(U)$ 一般并不连续，更糟的是它只在 $U$ 中几乎处处有定义。由于 $\partial U$ 具有 $n$ 维 Lebesgue 零测度，我们无法直接给“$u$ 限制在 $\partial U$ 上”这一表达式以任何意义。迹算子的概念解决了这一问题。

\begin{theorem}[$W^{k,p}$ 的迹与零迹刻画]\label{theorem:ch05-零迹函数-一般}
设 $k\geqslant1$，$1\leqslant p<\infty$，$U\subseteq\mathbb{R}^{n}$ 是有界开集，$\partial U$ 为 $C^{k}$。则存在有界线性算子
\begin{align*}
T_{k}:W^{k,p}(U)\longrightarrow\prod_{|\alpha|\leqslant k-1}L^{p}(\partial U),
\end{align*}
使得：
\begin{enumerate}[(i)]
\item 若 $u\in C^{k}(\overline{U})$，则
\begin{align*}
T_{k}u=\bigl(D^{\alpha}u|_{\partial U}\bigr)_{|\alpha|\leqslant k-1};
\end{align*}
\item 存在只依赖于 $n,k,p,U$ 的常数 $C$，使得对每个 $u\in W^{k,p}(U)$，
\begin{align}
\|T_{k}u\|_{\prod_{|\alpha|\leqslant k-1}L^{p}(\partial U)}
\leqslant C\|u\|_{W^{k,p}(U)};
\label{eq:ch05-一般零迹-迹估计}
\end{align}
\item
\begin{align}
W^{k,p}_{0}(U)=\left\{u\in W^{k,p}(U):T_{k}u=0\right\}.
\label{eq:ch05-一般零迹-核}
\end{align}
\end{enumerate}
\end{theorem}
\begin{proof}
\noindent$\mathbf{Step}\bm{1}$.{\heiti 在标准半空间上证明 $W^{1,p}$ 的迹估计.}

记 $\mathbb{R}^{n}_{+}=\{x=(x',x_n):x'\in\mathbb{R}^{n-1},\,x_n>0\}$，$\partial\mathbb{R}^{n}_{+}=\mathbb{R}^{n-1}=\{x_n=0\}$。

设 $g\in C^{1}(\overline{\mathbb{R}^{n}_{+}})\cap W^{1,p}(\mathbb{R}^{n}_{+})$，且 $\operatorname{spt}g$ 紧含于 $\overline{\mathbb{R}^{n}_{+}}$。对任意 $x=(x',x_n)\in\mathbb{R}^{n}_{+}$，由微积分基本定理，
\begin{align*}
g(x',0)=g(x',x_n)-\int_{0}^{x_n}\partial_n g(x',t)\,\mathrm{d}t.
\end{align*}
于是由 $(a+b)^{p}\leqslant2^{p-1}(a^{p}+b^{p})$ 及 Hölder 不等式，
\begin{align}
|g(x',0)|^{p}
&\leqslant2^{p-1}\left(|g(x',x_n)|^{p}+\left|\int_{0}^{x_n}\partial_n g(x',t)\,\mathrm{d}t\right|^{p}\right)
\notag\\
&\leqslant2^{p-1}\left(|g(x',x_n)|^{p}+x_n^{p-1}\int_{0}^{x_n}|\partial_n g(x',t)|^{p}\,\mathrm{d}t\right).
\label{eq:ch05-一般零迹-半空间逐点}
\end{align}
对 $x'\in\mathbb{R}^{n-1}$ 积分，得
\begin{align}
\int_{\mathbb{R}^{n-1}}|g(x',0)|^{p}\,\mathrm{d}x'
&\leqslant2^{p-1}\int_{\mathbb{R}^{n-1}}|g(x',x_n)|^{p}\,\mathrm{d}x'
\notag\\
&\qquad+2^{p-1}x_n^{p-1}\int_{0}^{x_n}\int_{\mathbb{R}^{n-1}}|\partial_n g(x',t)|^{p}\,\mathrm{d}x'\,\mathrm{d}t.
\label{eq:ch05-一般零迹-半空间积分}
\end{align}
将 \eqref{eq:ch05-一般零迹-半空间积分} 对 $x_n\in(0,1)$ 积分，得
\begin{align*}
\int_{\mathbb{R}^{n-1}}|g(x',0)|^{p}\,\mathrm{d}x'
&\leqslant2^{p-1}\int_{0}^{1}\int_{\mathbb{R}^{n-1}}|g(x',x_n)|^{p}\,\mathrm{d}x'\,\mathrm{d}x_n\\
&\qquad+2^{p-1}\int_{0}^{1}x_n^{p-1}\int_{0}^{x_n}\int_{\mathbb{R}^{n-1}}|\partial_n g(x',t)|^{p}\,\mathrm{d}x'\,\mathrm{d}t\,\mathrm{d}x_n.
\end{align*}
对第二个积分交换积分次序：
\begin{align*}
\int_{0}^{1}x_n^{p-1}\int_{0}^{x_n}|\partial_n g(x',t)|^{p}\,\mathrm{d}t\,\mathrm{d}x_n
&=\int_{0}^{1}\left(\int_{t}^{1}x_n^{p-1}\,\mathrm{d}x_n\right)|\partial_n g(x',t)|^{p}\,\mathrm{d}t\\
&\leqslant\frac{1}{p}\int_{0}^{1}|\partial_n g(x',t)|^{p}\,\mathrm{d}t.
\end{align*}
于是
\begin{align}
\int_{\mathbb{R}^{n-1}}|g(x',0)|^{p}\,\mathrm{d}x'
\leqslant C\left(\int_{\mathbb{R}^{n}_{+}}|g|^{p}\,\mathrm{d}x+\int_{\mathbb{R}^{n}_{+}}|Dg|^{p}\,\mathrm{d}x\right)
=C\|g\|_{W^{1,p}(\mathbb{R}^{n}_{+})}^{p},
\label{eq:ch05-一般零迹-半空间迹估计}
\end{align}
其中 $C$ 只依赖于 $p$。这就证明了半空间上的迹估计。

\noindent$\mathbf{Step}\bm{2}$.{\heiti 在 $C^{k}(\overline{U})$ 上定义 $T_{k}$ 并建立估计.}

由 $\partial U$ 紧及\cref{definition:appC-Ck边界}，存在有限个开集 $W_{i}$（$i=1,\cdots,N$）覆盖 $\partial U$，以及 $C^{k}$ 微分同胚
\begin{align*}
\Phi_{i}:W_{i}\to B(0,r_{i}),\qquad\Psi_{i}=\Phi_{i}^{-1},
\end{align*}
使得
\begin{align}
\Phi_{i}(W_{i}\cap U)=B^{+}(0,r_{i}),\qquad
\Phi_{i}(W_{i}\cap\partial U)=B(0,r_{i})\cap\{x_n=0\},
\label{eq:ch05-一般零迹-拉直}
\end{align}
其中 $B^{+}(0,r_{i})=B(0,r_{i})\cap\{x_n>0\}$。再取 $W_{0}\subset\subset U$ 使得
\begin{align}
\overline{U}\subseteq\bigcup_{i=0}^{N}W_{i}.
\label{eq:ch05-一般零迹-覆盖}
\end{align}
由\hyperref[corollary:紧集上的光滑单位分解定理]{紧集上的光滑单位分解定理}，存在 $\zeta_{i}\in C_{c}^{\infty}(W_{i})$，使得
\begin{align}
0\leqslant\zeta_{i}\leqslant1,\qquad
\sum_{i=0}^{N}\zeta_{i}(x)=1\quad(x\in\overline{U}).
\label{eq:ch05-一般零迹-单位分解}
\end{align}

设 $u\in C^{k}(\overline{U})$，$|\alpha|\leqslant k-1$。定义
\begin{align}
T_{k}u\triangleq\bigl(D^{\alpha}u|_{\partial U}\bigr)_{|\alpha|\leqslant k-1}.
\label{eq:ch05-一般零迹-Tk定义}
\end{align}
（这里 $D^{\alpha}u\in C^{1}(\overline{U})$，故 $D^{\alpha}u|_{\partial U}$ 是连续函数，从而属于 $L^{p}(\partial U)$。）

下面估计 $\|T_{k}u\|$。由于 $\sum_{i=0}^{N}\zeta_{i}=1$ 于 $\overline{U}$，有
\begin{align*}
D^{\alpha}u=\sum_{i=0}^{N}D^{\alpha}(\zeta_{i}u),
\end{align*}
故
\begin{align}
\|D^{\alpha}u\|_{L^{p}(\partial U)}^{p}
\leqslant C\sum_{i=0}^{N}\|D^{\alpha}(\zeta_{i}u)\|_{L^{p}(\partial U\cap W_{i})}^{p}.
\label{eq:ch05-一般零迹-分解估计}
\end{align}
对于 $i=0$，$\zeta_{0}u$ 的支集紧含于 $W_{0}\subset\subset U$，故 $D^{\alpha}(\zeta_{0}u)$ 在 $\partial U$ 的某邻域内为零，于是
\begin{align}
\|D^{\alpha}(\zeta_{0}u)\|_{L^{p}(\partial U\cap W_{0})}=0.
\label{eq:ch05-一般零迹-i0迹}
\end{align}

对于 $i\geqslant1$，令
\begin{align}
g_{i}(y)\triangleq D^{\alpha}(\zeta_{i}u)(\Psi_{i}(y)),\qquad y\in B^{+}(0,r_{i}).
\label{eq:ch05-一般零迹-gi}
\end{align}
由于 $\zeta_{i}u\in C^{k}(\overline{U})$，$|\alpha|\leqslant k-1$，故 $D^{\alpha}(\zeta_{i}u)\in C^{1}(\overline{U})$；又 $\Psi_{i}\in C^{k}$，故
\begin{align*}
g_{i}\in C^{1}(\overline{B^{+}(0,r_{i})}),\qquad
\operatorname{spt}g_{i}\ \text{紧含于}\ \overline{B^{+}(0,r_{i})}.
\end{align*}
对 $g_{i}$ 应用 $\mathbf{Step}\bm{1}$ 的半空间迹估计（把 $\mathbb{R}^{n}_{+}$ 换为 $B^{+}(0,r_{i})$，这只需在估计中再乘一个截断函数，不改变结论），得到
\begin{align}
\int_{B(0,r_{i})\cap\{y_n=0\}}|g_{i}(y',0)|^{p}\,\mathrm{d}y'
\leqslant C\|g_{i}\|_{W^{1,p}(B^{+}(0,r_{i}))}^{p}.
\label{eq:ch05-一般零迹-gi估计}
\end{align}
另一方面，变量替换 $x=\Psi_{i}(y)$ 给出
\begin{align}
\int_{B(0,r_{i})\cap\{y_n=0\}}|g_{i}(y',0)|^{p}\,\mathrm{d}y'
&=\int_{W_{i}\cap\partial U}|D^{\alpha}(\zeta_{i}u)(x)|^{p}J_{i}(x)\,\mathrm{d}S_{x},
\label{eq:ch05-一般零迹-曲面元}
\end{align}
其中 $J_{i}$ 是曲面测度的 Jacobi 因子。由于 $\partial U$ 紧且 $\Phi_{i}$ 是 $C^{1}$ 微分同胚，$J_{i}$ 在 $W_{i}\cap\partial U$ 上有正的上、下界。故存在 $C_{i}>0$，使得
\begin{align}
\|D^{\alpha}(\zeta_{i}u)\|_{L^{p}(\partial U\cap W_{i})}^{p}
\leqslant C_{i}\int_{B(0,r_{i})\cap\{y_n=0\}}|g_{i}(y',0)|^{p}\,\mathrm{d}y'.
\label{eq:ch05-一般零迹-曲面比较}
\end{align}
又由\cref{lemma:ch05-复合Wkp}（取 $k=1$），
\begin{align}
\|g_{i}\|_{W^{1,p}(B^{+}(0,r_{i}))}
\leqslant C\|D^{\alpha}(\zeta_{i}u)\|_{W^{1,p}(W_{i}\cap U)}
\leqslant C\|u\|_{W^{k,p}(U)}.
\label{eq:ch05-一般零迹-gi范数}
\end{align}
结合 \eqref{eq:ch05-一般零迹-gi估计}--\eqref{eq:ch05-一般零迹-gi范数}，
\begin{align}
\|D^{\alpha}(\zeta_{i}u)\|_{L^{p}(\partial U\cap W_{i})}
\leqslant C\|u\|_{W^{k,p}(U)}.
\label{eq:ch05-一般零迹-zi迹估计}
\end{align}
将 \eqref{eq:ch05-一般零迹-i0迹} 与 \eqref{eq:ch05-一般零迹-zi迹估计} 代入 \eqref{eq:ch05-一般零迹-分解估计}，并对 $|\alpha|\leqslant k-1$ 求和，得到
\begin{align}
\|T_{k}u\|_{\prod_{|\alpha|\leqslant k-1}L^{p}(\partial U)}
\leqslant C\|u\|_{W^{k,p}(U)},\qquad u\in C^{k}(\overline{U}).
\label{eq:ch05-一般零迹-Ck估计}
\end{align}

\noindent$\mathbf{Step}\bm{3}$.{\heiti 用 $C^{\infty}(\overline{U})$ 逼近把 $T_{k}$ 延拓到 $W^{k,p}(U)$.}

由\hyperref[theorem:ch05-直到边界光滑的函数整体逼近定理]{直到边界光滑的函数整体逼近定理}，$C^{\infty}(\overline{U})$ 在 $W^{k,p}(U)$ 中稠密。

设 $u\in W^{k,p}(U)$。取 $u_{m}\in C^{\infty}(\overline{U})$，使得
\begin{align*}
\|u_{m}-u\|_{W^{k,p}(U)}\longrightarrow0\quad(m\to\infty).
\end{align*}
由 \eqref{eq:ch05-一般零迹-Ck估计}，
\begin{align*}
\|T_{k}u_{m}-T_{k}u_{l}\|_{\prod L^{p}(\partial U)}
\leqslant C\|u_{m}-u_{l}\|_{W^{k,p}(U)}\longrightarrow0\quad(m,l\to\infty),
\end{align*}
故 $\{T_{k}u_{m}\}$ 是 $\prod_{|\alpha|\leqslant k-1}L^{p}(\partial U)$ 中的 Cauchy 列。由 $L^{p}(\partial U)$ 的完备性，存在极限，定义
\begin{align}
T_{k}u\triangleq\lim_{m\to\infty}T_{k}u_{m}.
\label{eq:ch05-一般零迹-Tk延拓}
\end{align}

下证 $T_{k}u$ 与逼近序列无关。设 $v_{m}\in C^{\infty}(\overline{U})$ 也满足 $\|v_{m}-u\|_{W^{k,p}(U)}\to0$。则
\begin{align*}
\|T_{k}u_{m}-T_{k}v_{m}\|_{\prod L^{p}(\partial U)}
\leqslant C\|u_{m}-v_{m}\|_{W^{k,p}(U)}
\leqslant C\left(\|u_{m}-u\|_{W^{k,p}(U)}+\|u-v_{m}\|_{W^{k,p}(U)}\right)\longrightarrow0,
\end{align*}
故两个极限相同。于是 $T_{k}$ 在 $W^{k,p}(U)$ 上良定义。

由 \eqref{eq:ch05-一般零迹-Ck估计} 取极限，得 \eqref{eq:ch05-一般零迹-迹估计}。由 $T_{k}$ 在 $C^{\infty}(\overline{U})$ 上的线性性与极限的线性性，$T_{k}$ 线性。若 $u\in C^{k}(\overline{U})$，取 $u_{m}=u$（或由 $C^{k}(\overline{U})$ 中的函数可用 $C^{\infty}(\overline{U})$ 函数在 $W^{k,p}$ 中逼近，且 $T_{k}$ 连续），即得 (i)。这就完成了迹算子 $T_{k}$ 的构造与性质 (i)(ii)。

\noindent$\mathbf{Step}\bm{4}$.{\heiti 证明半空间上带零迹函数的 $L^{p}$ 估计.}

设 $g\in W^{1,p}(\mathbb{R}^{n}_{+})$，$\operatorname{spt}g$ 紧含于 $\overline{\mathbb{R}^{n}_{+}}$，且其迹 $Tg=0$（这里 $T$ 是 $\mathbf{Step}\bm{1}$--$\bm{3}$ 的构造在 $k=1$、$U=\mathbb{R}^{n}_{+}$ 的半空间版本；由于 $\operatorname{spt}g$ 紧，可用 $\mathbf{Step}\bm{1}$ 的估计与 $C_{c}^{\infty}(\overline{\mathbb{R}^{n}_{+}})$ 逼近定义）。我们证明：对几乎处处的 $x_n>0$，
\begin{align}
\int_{\mathbb{R}^{n-1}}|g(x',x_n)|^{p}\,\mathrm{d}x'
\leqslant C x_n^{p-1}\int_{0}^{x_n}\int_{\mathbb{R}^{n-1}}|Dg(x',t)|^{p}\,\mathrm{d}x'\,\mathrm{d}t.
\label{eq:ch05-一般零迹-零迹估计}
\end{align}

取 $g_{m}\in C_{c}^{\infty}(\overline{\mathbb{R}^{n}_{+}})$，使得
\begin{align*}
g_{m}\longrightarrow g\quad\text{于}\ W^{1,p}(\mathbb{R}^{n}_{+}),\qquad
g_{m}(\cdot,0)\longrightarrow Tg=0\quad\text{于}\ L^{p}(\mathbb{R}^{n-1}).
\end{align*}
（存在性：先由 $C^{\infty}(\overline{\mathbb{R}^{n}_{+}})$ 在 $W^{1,p}$ 中稠密取逼近列，再乘以截断函数；迹的收敛由 $\mathbf{Step}\bm{1}$ 的估计得到。）对 $g_{m}$ 用 \eqref{eq:ch05-一般零迹-半空间逐点}：
\begin{align*}
|g_{m}(x',x_n)|^{p}
\leqslant2^{p-1}\left(|g_{m}(x',0)|^{p}+x_n^{p-1}\int_{0}^{x_n}|\partial_n g_{m}(x',t)|^{p}\,\mathrm{d}t\right).
\end{align*}
对 $x'$ 积分：
\begin{align*}
\int_{\mathbb{R}^{n-1}}|g_{m}(x',x_n)|^{p}\,\mathrm{d}x'
&\leqslant2^{p-1}\int_{\mathbb{R}^{n-1}}|g_{m}(x',0)|^{p}\,\mathrm{d}x'\\
&\qquad+2^{p-1}x_n^{p-1}\int_{0}^{x_n}\int_{\mathbb{R}^{n-1}}|\partial_n g_{m}(x',t)|^{p}\,\mathrm{d}x'\,\mathrm{d}t.
\end{align*}
令 $m\to\infty$。左端由 $g_{m}\to g$ 于 $L^{p}(\mathbb{R}^{n}_{+})$ 及\hyperref[Real Analysis-theorem:实变函数-Fubini定理]{Fubini定理}（取子列几乎处处收敛）收敛到左端；右端第一项由 $g_{m}(\cdot,0)\to0$ 于 $L^{p}(\mathbb{R}^{n-1})$ 趋于 $0$，第二项由 $\partial_n g_{m}\to\partial_n g$ 于 $L^{p}$ 收敛到相应的积分。于是对几乎处处的 $x_n>0$，\eqref{eq:ch05-一般零迹-零迹估计} 成立。

\noindent$\mathbf{Step}\bm{5}$.{\heiti 证明标准半空间上的零迹定理.}

设
\begin{align*}
v\in W^{k,p}(\mathbb{R}^{n}_{+}),\qquad
\operatorname{spt}v\ \text{紧含于}\ \overline{\mathbb{R}^{n}_{+}},
\end{align*}
且
\begin{align}
T(D^{\alpha}v)=0\quad\text{于}\ \partial\mathbb{R}^{n}_{+},\quad\forall\,|\alpha|\leqslant k-1.
\label{eq:ch05-一般零迹-半空间假设}
\end{align}
证明 $v\in W^{k,p}_{0}(\mathbb{R}^{n}_{+})$。

先证明高阶平坦性估计：对每个 $0\leqslant j\leqslant k-1$，每个 $|\alpha|=j$，以及几乎处处的 $x_n>0$，
\begin{align}
\int_{\mathbb{R}^{n-1}}|D^{\alpha}v(x',x_n)|^{p}\,\mathrm{d}x'
\leqslant C x_n^{(k-j)p-1}\int_{0}^{x_n}\int_{\mathbb{R}^{n-1}}|D^{k}v(x',s)|^{p}\,\mathrm{d}x'\,\mathrm{d}s.
\label{eq:ch05-一般零迹-平坦}
\end{align}
对 $j$ 作反向归纳。

当 $j=k-1$ 时，对 $|\alpha|=k-1$，$D^{\alpha}v\in W^{1,p}(\mathbb{R}^{n}_{+})$，$\operatorname{spt}(D^{\alpha}v)$ 紧含于 $\overline{\mathbb{R}^{n}_{+}}$，且 $T(D^{\alpha}v)=0$。对 $g=D^{\alpha}v$ 应用 \eqref{eq:ch05-一般零迹-零迹估计}：
\begin{align*}
\int_{\mathbb{R}^{n-1}}|D^{\alpha}v(x',x_n)|^{p}\,\mathrm{d}x'
\leqslant C x_n^{p-1}\int_{0}^{x_n}\int_{\mathbb{R}^{n-1}}|D(D^{\alpha}v)(x',s)|^{p}\,\mathrm{d}x'\,\mathrm{d}s.
\end{align*}
由于 $|D(D^{\alpha}v)|\leqslant|D^{k}v|$，即得 $j=k-1$ 的情形。

设 \eqref{eq:ch05-一般零迹-平坦} 对一切 $|\beta|=j+1$ 成立（$0\leqslant j\leqslant k-2$）。固定 $|\alpha|=j$。同上，$D^{\alpha}v\in W^{1,p}(\mathbb{R}^{n}_{+})$，$\operatorname{spt}(D^{\alpha}v)$ 紧含于 $\overline{\mathbb{R}^{n}_{+}}$，$T(D^{\alpha}v)=0$。由 \eqref{eq:ch05-一般零迹-零迹估计}，
\begin{align}
\int_{\mathbb{R}^{n-1}}|D^{\alpha}v(x',x_n)|^{p}\,\mathrm{d}x'
\leqslant C x_n^{p-1}\int_{0}^{x_n}\int_{\mathbb{R}^{n-1}}|D(D^{\alpha}v)(x',s)|^{p}\,\mathrm{d}x'\,\mathrm{d}s.
\label{eq:ch05-一般零迹-归纳1}
\end{align}
而
\begin{align*}
|D(D^{\alpha}v)(x',s)|^{p}\leqslant C\sum_{i=1}^{n}|D^{\alpha+e_i}v(x',s)|^{p}.
\end{align*}
对每个 $|\alpha+e_i|=j+1$，由归纳假设，
\begin{align}
\int_{\mathbb{R}^{n-1}}|D^{\alpha+e_i}v(x',s)|^{p}\,\mathrm{d}x'
\leqslant C s^{(k-j-1)p-1}\int_{0}^{s}\int_{\mathbb{R}^{n-1}}|D^{k}v(x',t)|^{p}\,\mathrm{d}x'\,\mathrm{d}t.
\label{eq:ch05-一般零迹-归纳2}
\end{align}
因为 $s\leqslant x_n$，有
\begin{align*}
\int_{0}^{s}\int_{\mathbb{R}^{n-1}}|D^{k}v|^{p}
\leqslant\int_{0}^{x_n}\int_{\mathbb{R}^{n-1}}|D^{k}v|^{p}.
\end{align*}
将 \eqref{eq:ch05-一般零迹-归纳2} 代入 \eqref{eq:ch05-一般零迹-归纳1}：
\begin{align*}
\int_{\mathbb{R}^{n-1}}|D^{\alpha}v(x',x_n)|^{p}\,\mathrm{d}x'
&\leqslant C x_n^{p-1}\int_{0}^{x_n}s^{(k-j-1)p-1}\,\mathrm{d}s
\int_{0}^{x_n}\int_{\mathbb{R}^{n-1}}|D^{k}v|^{p}\\
&=C' x_n^{(k-j)p-1}\int_{0}^{x_n}\int_{\mathbb{R}^{n-1}}|D^{k}v(x',t)|^{p}\,\mathrm{d}x'\,\mathrm{d}t.
\end{align*}
反向归纳完成。

下面作截断。取 $\zeta\in C^{\infty}(\mathbb{R})$，$0\leqslant\zeta\leqslant1$，$\zeta=1$ 于 $(-\infty,1]$，$\zeta=0$ 于 $[2,\infty)$。对 $\delta>0$ 令
\begin{align*}
\zeta_{\delta}(x)\triangleq\zeta\left(\frac{x_n}{\delta}\right),\qquad
w_{\delta}\triangleq v(1-\zeta_{\delta}).
\end{align*}
则 $\operatorname{spt}w_{\delta}$ 紧含于 $\{x_n\geqslant\delta\}$，且 $w_{\delta}\in W^{k,p}(\mathbb{R}^{n}_{+})$。由\hyperref[theorem:ch05-弱导数性质-4]{Leibniz 公式}，对 $|\alpha|\leqslant k$，
\begin{align}
D^{\alpha}w_{\delta}
=(1-\zeta_{\delta})D^{\alpha}v
-\sum_{\beta_n=1}^{|\alpha|}\binom{\alpha}{\beta_n e_n}\delta^{-\beta_n}\zeta^{(\beta_n)}\left(\frac{x_n}{\delta}\right)D^{\alpha-\beta_n e_n}v.
\label{eq:ch05-一般零迹-截断导数}
\end{align}
第一项由\hyperref[Real Analysis-theorem:控制收敛定理]{控制收敛定理}（控制函数 $|D^{\alpha}v|\in L^{p}$）收敛到 $D^{\alpha}v$。对求和中的每一项，用 \eqref{eq:ch05-一般零迹-平坦}（取 $j=|\alpha|-\beta_n$）：
\begin{align*}
\int_{\mathbb{R}^{n-1}}|D^{\alpha-\beta_n e_n}v(x',x_n)|^{p}\,\mathrm{d}x'
\leqslant C x_n^{(k-|\alpha|+\beta_n)p-1}\int_{0}^{x_n}\int_{\mathbb{R}^{n-1}}|D^{k}v|^{p}
\leqslant C x_n^{\beta_n p-1}\int_{0}^{2\delta}\int_{\mathbb{R}^{n-1}}|D^{k}v|^{p}
\end{align*}
（因为 $x_n\leqslant2\delta$ 且 $k-|\alpha|+\beta_n\geqslant\beta_n$）。于是该项的 $L^{p}$ 范数的 $p$ 次方不超过
\begin{align*}
C\delta^{-p\beta_n}\left(\int_{\delta}^{2\delta}x_n^{\beta_n p-1}\,\mathrm{d}x_n\right)
\int_{0}^{2\delta}\int_{\mathbb{R}^{n-1}}|D^{k}v|^{p}
\leqslant C\int_{0}^{2\delta}\int_{\mathbb{R}^{n-1}}|D^{k}v|^{p}\longrightarrow0
\end{align*}
（$\delta\to0^{+}$），因为 $D^{k}v\in L^{p}$ 而带状区域测度趋于零。故
\begin{align}
D^{\alpha}w_{\delta}\longrightarrow D^{\alpha}v\quad\text{于}\ L^{p}(\mathbb{R}^{n}_{+}),\qquad|\alpha|\leqslant k.
\label{eq:ch05-一般零迹-截断收敛}
\end{align}
又 $w_{\delta}\to v$ 于 $L^{p}$，故 $w_{\delta}\to v$ 于 $W^{k,p}(\mathbb{R}^{n}_{+})$。

最后磨光。固定 $\delta>0$，取标准磨光子 $\eta_{\varepsilon}$，$0<\varepsilon<\delta/2$，令
\begin{align*}
w_{\delta,\varepsilon}\triangleq\eta_{\varepsilon}*w_{\delta}.
\end{align*}
由\hyperref[theorem:appC-光滑算子的性质]{磨光子的性质}，$w_{\delta,\varepsilon}\in C^{\infty}(\mathbb{R}^{n})$，
\begin{align*}
\operatorname{spt}w_{\delta,\varepsilon}\subseteq\operatorname{spt}w_{\delta}+\overline{B(0,\varepsilon)}
\subseteq\{x_n\geqslant\delta-\varepsilon\}\subseteq\{x_n>0\},
\end{align*}
故 $w_{\delta,\varepsilon}\in C_{c}^{\infty}(\mathbb{R}^{n}_{+})$，且
\begin{align*}
w_{\delta,\varepsilon}\longrightarrow w_{\delta}\quad\text{于}\ W^{k,p}(\mathbb{R}^{n}_{+})\quad(\varepsilon\to0^{+}).
\end{align*}
取 $\delta_{m}\to0^{+}$、$\varepsilon_{m}<\delta_{m}/2$ 充分小，使得
\begin{align*}
\|w_{\delta_{m},\varepsilon_{m}}-v\|_{W^{k,p}(\mathbb{R}^{n}_{+})}\longrightarrow0.
\end{align*}
故 $v\in W^{k,p}_{0}(\mathbb{R}^{n}_{+})$。

\noindent$\mathbf{Step}\bm{6}$.{\heiti 用拉直与单位分解证明一般区域的零迹刻画.}

设 $u\in W^{k,p}(U)$ 满足
\begin{align}
T_{k}u=0.
\label{eq:ch05-一般零迹-假设}
\end{align}
证明 $u\in W^{k,p}_{0}(U)$。沿用 $\mathbf{Step}\bm{2}$ 的覆盖 $\{W_{i}\}_{i=0}^{N}$、拉直映射 $\Phi_{i},\Psi_{i}$ 与单位分解 $\{\zeta_{i}\}_{i=0}^{N}$。

对 $i=0$，$\zeta_{0}u$ 的支集紧含于 $W_{0}\subset\subset U$。取 $\delta>0$ 使 $\operatorname{spt}\zeta_{0}u+\overline{B(0,\delta)}\subset\subset U$，令 $(\zeta_{0}u)_{\varepsilon}=\eta_{\varepsilon}*(\zeta_{0}u)$，$0<\varepsilon<\delta$。则 $(\zeta_{0}u)_{\varepsilon}\in C_{c}^{\infty}(U)$，且由\hyperref[theorem:appC-光滑算子的性质]{磨光子的性质}知
\begin{align*}
(\zeta_{0}u)_{\varepsilon}\longrightarrow\zeta_{0}u\quad\text{于}\ W^{k,p}(U),
\end{align*}
故
\begin{align}
\zeta_{0}u\in W^{k,p}_{0}(U).
\label{eq:ch05-一般零迹-i0}
\end{align}

对 $i=1,\cdots,N$，令
\begin{align*}
v_{i}(y)\triangleq(\zeta_{i}u)(\Psi_{i}(y)),\qquad y\in B^{+}(0,r_{i}).
\end{align*}
由\cref{lemma:ch05-复合Wkp}，$v_{i}\in W^{k,p}(B^{+}(0,r_{i}))$，$\operatorname{spt}v_{i}$ 紧含于 $\overline{B^{+}(0,r_{i})}$。将 $v_{i}$ 零延拓到 $\mathbb{R}^{n}_{+}$，仍记为 $v_{i}$。则
\begin{align*}
v_{i}\in W^{k,p}(\mathbb{R}^{n}_{+}),\qquad
\operatorname{spt}v_{i}\ \text{紧含于}\ \overline{\mathbb{R}^{n}_{+}}.
\end{align*}

下面验证 $T(D^{\alpha}v_{i})=0$ 于 $\partial\mathbb{R}^{n}_{+}$。固定 $|\alpha|\leqslant k-1$。由\cref{lemma:ch05-复合Wkp} 的链式法则，在 $B^{+}(0,r_{i})$ 上
\begin{align}
D^{\alpha}v_{i}(y)
=\sum_{1\leqslant|\beta|\leqslant|\alpha|}(D^{\beta}(\zeta_{i}u))(\Psi_{i}(y))\,Q_{\alpha\beta}(y),
\label{eq:ch05-一般零迹-链式}
\end{align}
其中 $Q_{\alpha\beta}$ 是 $D^{\gamma}\Psi_{i}$（$1\leqslant|\gamma|\leqslant|\alpha|\leqslant k-1$）的多项式，故 $Q_{\alpha\beta}\in C^{0}(\overline{B^{+}(0,r_{i})})$。对 \eqref{eq:ch05-一般零迹-链式} 取迹。由于 $D^{\beta}(\zeta_{i}u)\in W^{1,p}(U)$ 有迹，而 $Q_{\alpha\beta}\circ\Phi_{i}$ 在 $\partial U\cap W_{i}$ 上连续，乘积的迹等于迹的乘积：
\begin{align}
T(D^{\alpha}v_{i})(y')
=\sum_{1\leqslant|\beta|\leqslant|\alpha|}T\bigl((D^{\beta}(\zeta_{i}u))\circ\Psi_{i}\bigr)(y')\,Q_{\alpha\beta}(y',0).
\label{eq:ch05-一般零迹-迹公式}
\end{align}
而
\begin{align*}
D^{\beta}(\zeta_{i}u)=\sum_{\gamma\leqslant\beta}\binom{\beta}{\gamma}D^{\gamma}\zeta_{i}\,D^{\beta-\gamma}u,
\end{align*}
其中 $D^{\gamma}\zeta_{i}$ 光滑，$D^{\beta-\gamma}u\in W^{1,p}(U)$，且 $|\beta-\gamma|\leqslant|\beta|\leqslant k-1$。由 \eqref{eq:ch05-一般零迹-假设}，
\begin{align*}
T(D^{\beta-\gamma}u)=0.
\end{align*}
由迹算子与光滑函数乘法的相容性（对光滑函数成立，再由 $T$ 的连续性与 $C^{\infty}(\overline U)$ 在 $W^{1,p}(U)$ 中的稠密性取极限），
\begin{align*}
T(D^{\beta}(\zeta_{i}u))
=\sum_{\gamma\leqslant\beta}\binom{\beta}{\gamma}(D^{\gamma}\zeta_{i})|_{\partial U}\,T(D^{\beta-\gamma}u)=0.
\end{align*}
于是
\begin{align*}
T\bigl((D^{\beta}(\zeta_{i}u))\circ\Psi_{i}\bigr)(y')
=T(D^{\beta}(\zeta_{i}u))(\Psi_{i}(y',0))=0.
\end{align*}
代入 \eqref{eq:ch05-一般零迹-迹公式}，得 $T(D^{\alpha}v_{i})=0$。由 $|\alpha|$ 的任意性，$v_{i}$ 满足 \eqref{eq:ch05-一般零迹-半空间假设}。由 $\mathbf{Step}\bm{5}$，$v_{i}\in W^{k,p}_{0}(\mathbb{R}^{n}_{+})$，从而
\begin{align}
v_{i}\in W^{k,p}_{0}(B^{+}(0,r_{i})).
\label{eq:ch05-一般零迹-vi}
\end{align}

由 \eqref{eq:ch05-一般零迹-vi}，存在 $v_{i,m}\in C_{c}^{\infty}(B^{+}(0,r_{i}))$，使得
\begin{align*}
v_{i,m}\longrightarrow v_{i}\quad\text{于}\ W^{k,p}(B^{+}(0,r_{i})).
\end{align*}
令 $u_{i,m}\triangleq v_{i,m}\circ\Phi_{i}$。则 $u_{i,m}\in C_{c}^{\infty}(W_{i}\cap U)$，且由\cref{lemma:ch05-复合Wkp}，
\begin{align*}
u_{i,m}\longrightarrow v_{i}\circ\Phi_{i}=\zeta_{i}u\quad\text{于}\ W^{k,p}(W_{i}\cap U).
\end{align*}
将 $u_{i,m}$ 与 $\zeta_{i}u$ 在 $U\setminus W_{i}$ 上零延拓，则 $u_{i,m}\in C_{c}^{\infty}(U)$ 且
\begin{align*}
u_{i,m}\longrightarrow\zeta_{i}u\quad\text{于}\ W^{k,p}(U).
\end{align*}
故
\begin{align}
\zeta_{i}u\in W^{k,p}_{0}(U),\qquad i=1,\cdots,N.
\label{eq:ch05-一般零迹-zi}
\end{align}

最后，由 \eqref{eq:ch05-一般零迹-单位分解}、\eqref{eq:ch05-一般零迹-i0}、\eqref{eq:ch05-一般零迹-zi} 及 $W^{k,p}_{0}(U)$ 的线性性，
\begin{align*}
u=\zeta_{0}u+\sum_{i=1}^{N}\zeta_{i}u\in W^{k,p}_{0}(U).
\end{align*}
这就证明了 \eqref{eq:ch05-一般零迹-核} 中“$\supseteq$”的包含关系。

反之，若 $u\in W^{k,p}_{0}(U)$，则由 $C_{c}^{\infty}(U)$ 在 $W^{k,p}(U)$ 中的稠密性，取 $u_{m}\in C_{c}^{\infty}(U)$，$u_{m}\to u$ 于 $W^{k,p}(U)$。对每个 $|\alpha|\leqslant k-1$，$D^{\alpha}u_{m}\in C_{c}^{\infty}(U)$，故 $T(D^{\alpha}u_{m})=0$。由 $T_{k}$ 的连续性，
\begin{align*}
T_{k}u=\lim_{m\to\infty}T_{k}u_{m}=0.
\end{align*}
故 $W^{k,p}_{0}(U)\subseteq\ker T_{k}$。结合上面的反向包含，\eqref{eq:ch05-一般零迹-核} 成立。
\end{proof}




\end{document}
